\documentclass[10pt,a4paper]{amsart}
\usepackage{amsmath,amssymb,amsthm,graphicx}
% Retained minimal preamble for the Stacks Project chapter more-algebra.tex (More on Algebra).
% The chapter file pinned-original-source.tex carries only \input{preamble} on its first line, so no
% preamble text lives inside the chapter itself.  The Stacks Project uses one single shared preamble for
% every chapter; the snapshot stages it identically as sources/stacks-chow/preamble.tex and
% sources/stacks-etale-cohomology/preamble.tex.  The block below is copied mechanically, line for line,
% from sources/stacks-chow/preamble.tex lines 14-18 (the Xy-pic setup) and 154-217 (the native theorem
% declarations together with the mathematical macros).  Nothing was added, removed or redefined; the
% wrapper class and the amsmath/amssymb/amsthm/graphicx base packages are supplied by the assembler.
% Chapter: more-algebra.tex -- More on Algebra.
% For commutative diagrams we use Xy-pic
\usepackage[all]{xy}

% We use 2cell for 2-commutative diagrams.
\xyoption{2cell}

% Theorem environments.
%
\theoremstyle{plain}
\newtheorem{theorem}[subsection]{Theorem}
\newtheorem{proposition}[subsection]{Proposition}
\newtheorem{lemma}[subsection]{Lemma}

\theoremstyle{definition}
\newtheorem{definition}[subsection]{Definition}
\newtheorem{example}[subsection]{Example}
\newtheorem{exercise}[subsection]{Exercise}
\newtheorem{situation}[subsection]{Situation}

\theoremstyle{remark}
\newtheorem{remark}[subsection]{Remark}
\newtheorem{remarks}[subsection]{Remarks}

\numberwithin{equation}{subsection}

% Macros
%
\def\lim{\mathop{\mathrm{lim}}\nolimits}
\def\colim{\mathop{\mathrm{colim}}\nolimits}
\def\Spec{\mathop{\mathrm{Spec}}}
\def\Hom{\mathop{\mathrm{Hom}}\nolimits}
\def\Ext{\mathop{\mathrm{Ext}}\nolimits}
\def\SheafHom{\mathop{\mathcal{H}\!\mathit{om}}\nolimits}
\def\SheafExt{\mathop{\mathcal{E}\!\mathit{xt}}\nolimits}
\def\Sch{\mathit{Sch}}
\def\Mor{\mathop{\mathrm{Mor}}\nolimits}
\def\Ob{\mathop{\mathrm{Ob}}\nolimits}
\def\Sh{\mathop{\mathit{Sh}}\nolimits}
\def\NL{\mathop{N\!L}\nolimits}
\def\CH{\mathop{\mathrm{CH}}\nolimits}
\def\proetale{{pro\text{-}\acute{e}tale}}
\def\etale{{\acute{e}tale}}
\def\QCoh{\mathit{QCoh}}
\def\Ker{\mathop{\mathrm{Ker}}}
\def\Im{\mathop{\mathrm{Im}}}
\def\Coker{\mathop{\mathrm{Coker}}}
\def\Coim{\mathop{\mathrm{Coim}}}

% Boxtimes
%
\DeclareMathSymbol{\boxtimes}{\mathbin}{AMSa}{"02}

%
% Macros for moduli stacks/spaces
%
\def\QCohstack{\mathcal{QC}\!\mathit{oh}}
\def\Cohstack{\mathcal{C}\!\mathit{oh}}
\def\Spacesstack{\mathcal{S}\!\mathit{paces}}
\def\Quotfunctor{\mathrm{Quot}}
\def\Hilbfunctor{\mathrm{Hilb}}
\def\Curvesstack{\mathcal{C}\!\mathit{urves}}
\def\Polarizedstack{\mathcal{P}\!\mathit{olarized}}
\def\Complexesstack{\mathcal{C}\!\mathit{omplexes}}
% \Pic is the operator that assigns to X its picard group, usage \Pic(X)
% \Picardstack_{X/B} denotes the Picard stack of X over B
% \Picardfunctor_{X/B} denotes the Picard functor of X over B
\def\Pic{\mathop{\mathrm{Pic}}\nolimits}
\def\Picardstack{\mathcal{P}\!\mathit{ic}}
\def\Picardfunctor{\mathrm{Pic}}
\def\Deformationcategory{\mathcal{D}\!\mathit{ef}}
\usepackage{cleveref}
\crefname{originallemma}{original lemma}{original lemmas}
\Crefname{originallemma}{Original lemma}{Original lemmas}
\makeatletter
\newlabel{more-algebra-lemma-injective-abelian}{{01D7}{source}{}{sourcepointer-1}{}}
\newlabel{more-algebra-lemma-injective-abelian@cref}{{[originallemma][0][]01D7}{[1][0][]source}{}{}{}}
\makeatother
\title{An abelian group is injective exactly when it is divisible}
\author{The Stacks Project Authors}
\date{stacks-project:more-algebra.tex:local-snapshot-20261012; GNU FDL 1.2 or later}
\begin{document}
\maketitle
\setcounter{subsection}{1}
\noindent\textit{Editorial note.} All delivered mathematical text is the author's own, copied from the pinned Stacks Project chapter file more-algebra.tex; no statement or proof step was rewritten, reordered or removed.\par\medskip
\noindent\textit{Editorial note.} The chapter file pinned-original-source.tex carries only \textbackslash{}input\{preamble\} on its first line, so no preamble text lives inside the chapter. The retained file sources/stacks-more-algebra/retained-preamble.tex is a mechanical, line-for-line copy of the project's single shared preamble as staged in the snapshot, namely sources/stacks-chow/preamble.tex lines 14-18 (Xy-pic setup) and 154-217 (theorem declarations and mathematical macros). Nothing was added, removed or redefined, no package was added, and the excerpt serves only as preamble, never as a span coordinate.\par\medskip
\noindent\textit{Editorial note.} Both counted spans start at the \textbackslash{}begin line of their environment, so the statement and the author's proof are delivered whole; surrounding prose of the section is deliberately outside the counted spans and is carried as standard prerequisites instead.\par\medskip
\noindent\textit{Editorial note.} The citation fields (source\_id, source\_version, authors, source\_url, license.name, source\_title) are identical to those of the cards already published from this chapter, so that the batch selector treats them as one source instead of deferring them as a same-source clash.\par\medskip
\noindent\textit{Allowed standard prerequisites (editorial).}\begin{itemize}
\item Injective object in an abelian category, here tested by the extension property for inclusions of subgroups. definition and discussion earlier in the chapter, source lines 13480-13551, immediately before the statement. Supplies the property that is being characterised and the form in which the necessity is proved.
\item Divisibility of an abelian group. recalled at source lines 13548-13551 just before the statement. Supplies the hypothesis in the sufficiency half and the element z used in the second case of the extension step.
\item Zorn's lemma applied to partial extensions, with glueing along a chain. used at proof lines 13575-13581. Produces a maximal partial extension, so that the whole problem reduces to the single step of extending across one extra element.
\end{itemize}\medskip
\ifdefined\hypertarget\hypertarget{sourcepointer-1}{}\fi
\noindent\textit{Source reference.} Original lemma 01D7: focused-seeds/AGP/sources/stacks-snapshot-20261011/tags.txt, source lines 1788--1788. The lemma's own canonical tag, quoted so that readers can resolve the label inside the delivered text; it points at the counted statement itself and carries no proof content.\par\medskip
\begin{lemma}
\label{lemma-injective-abelian}
An abelian group $J$ is an injective object in
the category of abelian groups if and only if $J$
is divisible.
\end{lemma}

\begin{proof}
Suppose that $J$ is not divisible. Then there exists
an $x \in J$ and $n \in \mathbf{N}$ such that there
is no $y \in J$ with $n y = x$. Then the morphism
$\mathbf{Z} \to J$, $m \mapsto mx$ does not extend
to $\frac{1}{n}\mathbf{Z} \supset \mathbf{Z}$. Hence
$J$ is not injective.

\medskip\noindent
Let $A \subset B$ be abelian groups.
Assume that $J$ is a divisible abelian group.
Let $\varphi : A \to J$ be a morphism.
Consider the set of homomorphisms $\varphi' : A' \to J$
with $A \subset A' \subset B$ and $\varphi'|_A = \varphi$.
Define $(A', \varphi') \geq (A'', \varphi'')$ if
and only if $A' \supset A''$ and $\varphi'|_{A''} = \varphi''$.
If $(A_i, \varphi_i)_{i \in I}$ is a totally
ordered collection of such pairs, then we obtain a map
$\bigcup_{i \in I} A_i \to J$ defined by $a \in A_i$
maps to $\varphi_i(a)$. Thus Zorn's lemma applies.
To conclude we have to show that if the pair
$(A', \varphi')$ is maximal then $A' = B$.
In other words, it suffices to show, given
any subgroup $A \subset B$, $A \not = B$ and
any $\varphi : A \to J$, then we can find
$\varphi' : A' \to J$ with $A \subset A' \subset B$
such that (a) the inclusion $A \subset A'$ is strict, and
(b) the morphism $\varphi'$ extends $\varphi$.

\medskip\noindent
To prove this, pick $x \in B$, $x \not \in A$.
If there exists no $n\in \mathbf{N}$ such that
$nx \in A$, then $A \oplus \mathbf{Z} \cong A + \mathbf{Z}x$.
Hence we can extend $\varphi$ to $A' = A + \mathbf{Z}x$
by using $\varphi$ on $A$ and mapping $x$ to zero for example.
If there does exist an $n \in \mathbf{N}$ such that
$nx \in A$, then let $n$ be the minimal such integer.
Let $z \in J$ be an element such that $nz = \varphi(nx)$.
Define a morphism $\tilde\varphi : A \oplus \mathbf{Z} \to J$ by
$(a, m) \mapsto \varphi(a) + mz$. By our choice of
$z$ the kernel of $\tilde \varphi$ contains the kernel
of the map $A \oplus \mathbf{Z} \to B$,
$(a, m) \mapsto a + mx$. Hence $\tilde \varphi$ factors
through the image $A' = A + \mathbf{Z}x$, and this extends the morphism
$\varphi$.
\end{proof}

\end{document}
